% --- Core Mathematical Packages ---
\usepackage{amsmath}
\usepackage{amssymb}
\usepackage{amsthm}
\usepackage{bm}

% --- Page Layout & Headers ---
\usepackage{geometry}
\geometry{
    a4paper,
    margin=1in,
    headheight=15pt
}
\usepackage{fancyhdr}
\pagestyle{fancy}
\fancyhf{}
\fancyfoot[C]{\thepage}
\lhead{\emph{\nouppercase{\leftmark}}}

% --- Typography, Lists & Helpers ---
\usepackage{enumitem}
\usepackage{etoolbox}
\usepackage{caption}
\usepackage{pgffor}
\usepackage{xcolor}
\usepackage{tikz}
\usetikzlibrary{arrows.meta, positioning}

% List formatting
\renewcommand{\theenumi}{\alph{enumi}}
\renewcommand{\labelenumi}{(\theenumi)}
\setlist[itemize,2]{label=$\circ$}
\setlist{
    topsep=3pt,
    itemsep=2pt,
    parsep=0pt,
    before=\leavevmode
}

% Paragraph layout: zero indentation, clean inter-paragraph glue
\setlength{\parindent}{0pt}
\setlength{\parskip}{5pt plus 2pt minus 1pt}

% Document Metadata
\author{ioshift}
\date{\today}
\title{\Huge\bfseries\ntitle}


%%%%%%%%%%%%%%%%%%%%%%%%%%%%%%%%%%%%%%%%%%%%%%%%%%%%%%%%%%%%%%%%%%%%%%%%%%%%%%%
%% Theorem Environments & Styling
%%%%%%%%%%%%%%%%%%%%%%%%%%%%%%%%%%%%%%%%%%%%%%%%%%%%%%%%%%%%%%%%%%%%%%%%%%%%%%%

\numberwithin{equation}{subsection}

% Custom Theorem Styles with built-in vertical glue
\newtheoremstyle{notestyle_thm}
  {6pt plus 2pt minus 1pt}  % Space above
  {6pt plus 2pt minus 1pt}  % Space below
  {\normalfont}             % Body font
  {}                        % Indent amount
  {\bfseries}               % Head font
  {.}                       % Punctuation after head
  { }                       % Space after head
  {}                        % Head spec

\newtheoremstyle{notestyle_remark}
  {5pt plus 2pt minus 1pt}  % Space above
  {5pt plus 2pt minus 1pt}  % Space below
  {\normalfont}             % Body font
  {}                        % Indent amount
  {\bfseries}               % Head font
  {.}                       % Punctuation after head
  { }                       % Space after head
  {}                        % Head spec

% Numbered theorem-like environments
\theoremstyle{notestyle_thm}
\newtheorem{thm}{Theorem}[subsection]
\renewcommand{\thethm}{%
  \ifnum\value{subsection}=0
    \thesection.\arabic{thm}%
  \else
    \thesubsection.\arabic{thm}%
  \fi
}

\newtheorem{axiom}[thm]{Axiom}
\newtheorem{conjecture}[thm]{Conjecture}
\newtheorem{cor}[thm]{Corollary}
\newtheorem{eg}[thm]{Example}
\newtheorem{prop}[thm]{Proposition}
\newtheorem{lemma}[thm]{Lemma}
\newtheorem{ex}[thm]{Exercise}
\newtheorem{definition}[thm]{Definition}
\newtheorem{notations}[thm]{Notation}

% Unnumbered / Starred environments
\theoremstyle{notestyle_remark}
\newtheorem*{aim}{Aim}
\newtheorem*{assumption}{Assumption}
\newtheorem*{claim}{Claim}
\newtheorem*{Defi}{Definition}
\newtheorem*{fact}{Fact}
\newtheorem*{law}{Law}
\newtheorem*{remark}{Remark}
\newtheorem*{note}{Note}
\newtheorem*{rrule}{Rule}
\newtheorem*{Prop}{Proposition}
\newtheorem*{question}{Question}
\newtheorem*{warning}{Warning}

% Proof & Solution Environments
\renewcommand{\proofname}{\textbf{Proof}}
\newenvironment{sol}{\begin{proof}[\textbf{Solution}]}{\end{proof}}
\newenvironment{counter}{\begin{proof}[\textbf{Counterexample}]}{\end{proof}}


%%%%%%%%%%%%%%%%%%%%%%%%%%%%%%%%%%%%%%%%%%%%%%%%%%%%%%%%%%%%%%%%%%%%%%%%%%%%%%%
%% Mathematical Macros & Symbols
%%%%%%%%%%%%%%%%%%%%%%%%%%%%%%%%%%%%%%%%%%%%%%%%%%%%%%%%%%%%%%%%%%%%%%%%%%%%%%%

% Standard blackboard-bold sets
\newcommand{\C}{\mathbb{C}}
\newcommand{\R}{\mathbb{R}}
\newcommand{\Q}{\mathbb{Q}}
\newcommand{\Z}{\mathbb{Z}}
\newcommand{\N}{\mathbb{N}}
\newcommand{\F}{\mathbb{F}}

% Calligraphic collections
\newcommand{\topo}{\mathcal{T}}

% Mathematical Operators
\DeclareMathOperator{\re}{Re}
\DeclareMathOperator{\im}{Im}
\DeclareMathOperator{\interior}{Int}
\DeclareMathOperator{\lcm}{lcm}
\DeclareMathOperator{\diam}{diam}
\DeclareMathOperator{\id}{id}
\DeclareMathOperator{\Bd}{Bd}
\DeclareMathOperator{\card}{card}
\DeclareMathOperator{\supp}{supp}
\DeclareMathOperator{\image}{im}

% Delimiters & Accents
\newcommand{\abs}[1]{\left\lvert #1 \right\rvert}
\newcommand{\norm}[1]{\left\lVert #1 \right\rVert}
\newcommand{\floor}[1]{\left\lfloor #1 \right\rfloor}
\newcommand{\ceil}[1]{\left\lceil #1 \right\rceil}
\newcommand{\set}[1]{\left\{ #1 \right\}}
\newcommand{\closure}[2][3]{{}\mkern#1mu\overline{\mkern-#1mu#2}}
\newcommand{\group}[1]{\langle #1 \rangle}

% Symbols & Relations
\newcommand{\noteq}{\neq}
\newcommand{\nexist}{\nexists}
\newcommand{\aand}{\quad\text{and}\quad}
\newcommand{\res}[1]{\!\mid_{#1}}
\newcommand{\normal}{\unlhd}
\renewcommand{\tt}[1]{\quad\text{#1}\quad}
\renewcommand{\d}{\mathop{}\!\mathrm{d}}
\renewcommand{\emptyset}{\varnothing}

% Darboux / Riemann integrals
\newcommand{\upint}{\mathchoice%
    {\mkern13mu\overline{\vphantom{\intop}\mkern7mu}\mkern-20mu}%
    {\mkern7mu\overline{\vphantom{\intop}\mkern7mu}\mkern-14mu}%
    {\mkern7mu\overline{\vphantom{\intop}\mkern7mu}\mkern-14mu}%
    {\mkern7mu\overline{\vphantom{\intop}\mkern7mu}\mkern-14mu}%
  \int}
\newcommand{\lowint}{\mkern3mu\underline{\vphantom{\intop}\mkern7mu}\mkern-10mu\int}


%%%%%%%%%%%%%%%%%%%%%%%%%%%%%%%%%%%%%%%%%%%%%%%%%%%%%%%%%%%%%%%%%%%%%%%%%%%%%%%
%% Automatic Definition & Concept Glossary Engine
%%%%%%%%%%%%%%%%%%%%%%%%%%%%%%%%%%%%%%%%%%%%%%%%%%%%%%%%%%%%%%%%%%%%%%%%%%%%%%%

\newcounter{cnt}

\newcommand{\newterm}[3]{%
    \expandafter\gdef\csname term#1\endcsname{#2}%
    \expandafter\gdef\csname cat#1\endcsname{#3}%
}

% Definition environment with automatic glossary registration
\newenvironment{defi}[1]{%
    \refstepcounter{cnt}%
    \phantomsection\label{glo:\arabic{cnt}}%
    \newterm{\arabic{cnt}}{#1}{0}%
    \begin{definition}\textbf{#1}
    
}{%
    \end{definition}%
}

% Notation environment with automatic glossary registration
\newenvironment{notation}[1]{%
    \refstepcounter{cnt}%
    \phantomsection\label{glo:\arabic{cnt}}%
    \newterm{\arabic{cnt}}{#1}{0}%
    \begin{notations}\textbf{#1}
    
}{%
    \end{notations}%
}

% Ad-hoc glossary elements and category headers
\newcommand{\addelement}[1]{%
    \refstepcounter{cnt}%
    \phantomsection\label{glo:\arabic{cnt}}%
    \newterm{\arabic{cnt}}{#1}{0}%
}

\newcommand{\addclass}[1]{%
    \refstepcounter{cnt}%
    \phantomsection\label{glo:\arabic{cnt}}%
    \newterm{\arabic{cnt}}{#1}{1}%
}

% Glossary rendering command
\newcommand{\printglossary}{%
    \newpage
    \section{Glossary}
    \ifnum\value{cnt}>0
      \foreach \x in {1,...,\arabic{cnt}} {%
          \ifnum\csname cat\x\endcsname=0
              \noindent\hyperref[glo:\x]{\csname term\x\endcsname}\par\addvspace{0.4em}%
          \else
              \ifnum\x>1
                  \addvspace{1.2em}%
              \fi
              \noindent{\large\textbf{\csname term\x\endcsname}}\par\addvspace{0.6em}%
          \fi
      }%
    \fi
}


%%%%%%%%%%%%%%%%%%%%%%%%%%%%%%%%%%%%%%%%%%%%%%%%%%%%%%%%%%%%%%%%%%%%%%%%%%%%%%%
%% Hyperref (Loaded Last for Guaranteed Compatibility)
%%%%%%%%%%%%%%%%%%%%%%%%%%%%%%%%%%%%%%%%%%%%%%%%%%%%%%%%%%%%%%%%%%%%%%%%%%%%%%%

\usepackage[colorlinks=true, linkcolor=red, citecolor=red, urlcolor=blue]{hyperref}